\subsubsection{Pressure and Large Deviations of Weighted Kernels: Auditable Interface}\label{subsubsec:sync-kernel-weighted-pressure-ldp}
Let $\phi$ be an observable on the edges of the synchronization graph (e.g., the indicator function of output bit $e$, or the indicator "reaching a state containing $\overline{1}$"). Define the weighted matrix
\[
(B_\theta)_{ij}:=\sum_{e:i\to j}\exp\bigl(\theta\,\phi(e)\bigr),
\]
where the sum is over all edges from $i$ to $j$. Write
\[
\lambda(\theta):=\rho(B_\theta),\qquad P(\theta):=\log\lambda(\theta).
\]
Then $P(\theta)$ is the pressure function for the finite-type case, and $P(0)=\log 3$. Let
\[
I(\alpha):=\sup_{\theta\in\RR}\bigl(\theta\alpha-(P(\theta)-P(0))\bigr),
\]
then $I(\alpha)$ gives the large deviation rate function for the long-run average of $\phi$ (standard thermodynamic formalism; see \cite{Ruelle1978Thermodynamic,ParryPollicott1990Zeta}).

The weighted $\zeta$ is still given by
\[
\zeta_{B,\theta}(z)=\frac{1}{\det(I-zB_\theta)}
\]
and one can define the weighted primitive orbit spectrum
\[
a_n(\theta):=\mathrm{Tr}(B_\theta^n),\qquad
p_n(\theta)=\frac{1}{n}\sum_{d\mid n}\mu(d)\,a_{n/d}(d\theta),
\]
as well as the corresponding weighted Mertens constant family (obtained by taking the Abel finite part as $z\to 1/\lambda(\theta)$).
Auditable numerics are output by script \texttt{scripts/exp\_sync\_kernel\_weighted\_pressure.py}.

\subsubsection{Witt Decomposition of Weighted Primitive--Periodic}\label{app:witt-weighted}
Below we take output potential $\varphi(e)=e$, and let
\[
u=e^\theta,\qquad B(u)=B_0+uB_1,\qquad \Delta(z,u):=\det(I-zB(u)).
\]
Here $B_0$ (resp.\ $B_1$) counts only edges with output bit $e=0$ (resp.\ $e=1$).
Then
\[
\zeta(z,u)=\exp\!\left(\sum_{n\ge 1}\frac{\mathrm{Tr}(B(u)^n)}{n}z^n\right)=\Delta(z,u)^{-1}.
\]
Let
\[
a_n(u):=\mathrm{Tr}(B(u)^n),\qquad
P(z,u):=\sum_{n\ge 1}p_n(u)z^n,
\]
where
\(
p_n(u)=\sum_{\tau\ \mathrm{primitive},\,|\tau|=n}u^{S_\tau\varphi}
\)
is the weighted spectral polynomial of primitive orbits. The Euler product gives
\[
\zeta(z,u)=\prod_{\tau\ \mathrm{primitive}}\bigl(1-z^{|\tau|}u^{S_\tau\varphi}\bigr)^{-1}.
\]
Taking logarithm and expanding power repetition yields the Witt decomposition
\[
\log\zeta(z,u)=\sum_{m\ge 1}\frac{1}{m}P(z^m,u^m).
\]
By Möbius inversion we obtain
\[
P(z,u)=\sum_{m\ge 1}\frac{\mu(m)}{m}\log\zeta(z^m,u^m),
\]
comparing coefficients immediately gives
\[
p_n(u)=\frac{1}{n}\sum_{d\mid n}\mu(d)\,a_{n/d}(u^d)
\quad\Longleftrightarrow\quad
p_n(\theta)=\frac{1}{n}\sum_{d\mid n}\mu(d)\,a_{n/d}(d\theta).
\]

\subsubsection{Big-Witt/Necklace Bridge Package and Verifiable Congruences}\label{app:witt-bridge-package}
The above formula is not merely a "computable formula," but rather a standard coordinate transformation on Big-Witt (necklace ring): $a_n(u)$ are ghost components, $p_n(u)$ are Witt coordinates, and the Euler product is the plethystic exponential. Equivalently, for any parametrized matrix family $T(u)$, there is a strict identity
\[
\boxed{\ \zeta_T(z,u)=\mathrm{Exp}_W\bigl(\mathrm{Mob}_W(\mathrm{TrSeq}(T))\bigr)\ }
\]
where $\mathrm{TrSeq}(T)=(a_n(u))_{n\ge 1}$, $\mathrm{Mob}_W(a)=p$ gives $p_n(u)=\frac1n\sum_{d\mid n}\mu(d)\,a_{n/d}(u^d)$, and $\mathrm{Exp}_W$ is defined by
\(
\log\zeta=\sum_{m\ge 1}\frac{1}{m}P(z^m,u^m)
\)
. This "bridge package" directly produces a batch of verifiable arithmetic constraints:

\begin{proposition}[Functoriality of the commutative diagram: naturality under intertwiner]\label{prop:witt-bridge-functorial}
Take coefficient ring \(R\) (e.g.\ \(R=\ZZ[u]\) or \(R=\ZZ[u,u^{-1}]\)), and let objects in \(\mathbf{End}_R\) be triples \((V,T)\) (finitely generated free \(R\)-module with an \(R\)-linear endomorphism), with morphisms being intertwiners \(f:(V,T)\to (W,S)\) satisfying \(f\circ T=S\circ f\).
Let \(\mathbf{Seq}_R\) be the category of \(R\)-valued sequences (with component-wise maps), and define
\[
\mathsf{Ghost}:\mathbf{End}_R\to \mathbf{Seq}_R,\qquad
(V,T)\mapsto (a_n)_{n\ge 1},\ a_n:=\Tr(T^n).
\]
Then \(\mathsf{Ghost}\) is a functor, and for any intertwiner \(f\) we have \(a_n(T)=a_n(S)\) (hence \(\zeta_T=\zeta_S\)).
Now let \(\mathsf{Mob}\) be the Witt--Möbius transform with Adams operation \(\psi^d(u)=u^d\)
\[
\mathsf{Mob}(a)=p,\qquad
p_n(u):=\frac1n\sum_{d\mid n}\mu(d)\,a_{n/d}(u^d),
\]
and let \(\mathsf{Euler}\) be the construction that generates the Euler product from \(p\) (equivalently the plethystic exponential).
Then for each \((V,T)\) there is a strict identity
\[
\boxed{\;
\mathsf{Zeta}(V,T)
=\mathsf{Euler}\bigl(\mathsf{Mob}(\mathsf{Ghost}(V,T))\bigr)\; },
\qquad
\mathsf{Zeta}(V,T):=\exp\!\left(\sum_{n\ge 1}\frac{\Tr(T^n)}{n}z^n\right)=\det(I-zT)^{-1}.
\]
Thus the commutative diagram "trace (time ghost) \(\leftrightarrow\) primitive (space atom)" is not merely an analogy, but a strict natural transformation given by the Big-Witt coordinate transformation.
\end{proposition}

\begin{theorem}[Witt--Möbius Transform of Dirichlet Character Twist]\label{thm:witt-dirichlet-twist}
For any Dirichlet character \(\chi\), define the twisted \(\mathsf{Mob}\) transform
\[
p_n^{(\chi)}:=\frac{1}{n}\sum_{d\mid n}\mu(d)\,\chi(d)\,\psi^d\!\bigl(a_{n/d}\bigr)\in R,
\]
and define the twisted zeta generating function
\[
\mathsf{Zeta}_{\chi}(V,T)
:=\exp\!\Bigl(\sum_{n\ge 1}\frac{\chi(n)}{n}\,\Tr(T^{n})\,z^{n}\Bigr).
\]
Then there exists a strict identity
\[
\boxed{\;\mathsf{Zeta}_{\chi}(V,T)=\mathsf{Euler}\bigl(p^{(\chi)}\bigr)\; },
\]
and its "primitive stripping" (i.e.\ \(\mathsf{Mob}_{\chi}^{-1}\)) has Dirichlet coefficients \(\mu(n)\chi(n)\), completely isomorphic to the untwisted case \(\mathcal{P}(s)=\sum_{m\ge 1}\frac{\mu(m)}{m}\log\mathcal{Z}(ms)\).
\end{theorem}
\begin{proof}
By the twisted ghost components \(a_n^{(\chi)}:=\chi(n)\Tr(T^n)\) and the standard Witt/Möbius inversion \(n\,p_n=\sum_{d\mid n}\mu(d)\,\psi^d(a_{n/d})\), substituting \(a_{n/d}^{(\chi)}=\chi(n/d)\Tr(T^{n/d})\) gives \(p_n^{(\chi)}=\frac{1}{n}\sum_{d\mid n}\mu(d)\chi(d)\,\psi^d(a_{n/d})\); the Euler product is given by the correspondence between \(\mathsf{Zeta}_{\chi}=\exp\bigl(\sum_n \chi(n)\Tr(T^n)z^n/n\bigr)\) and the primitive expansion \(\log\mathsf{Zeta}_{\chi}=\sum_m p_m^{(\chi)}(z^m)/m\). The primitive stripping Möbius inversion kernel is \(\mu(n)\chi(n)\). QED.
\end{proof}
This theorem provides a characterization of the internal orthogonal representation decomposition of the system: when phase compression introduces character channels (e.g., branch memory modulo \(4\)), characters are not external labels, but rather necessary twists of the Big-Witt natural transformation on the Adams structure; the channelization decomposition of the phase layer is completely characterized by the representation-theoretic basis vectors given by Dirichlet characters.

\begin{definition}[Adams--Dirichlet Renormalization Gate (Arithmetic Rescaling)]\label{def:witt-adams-dirichlet-renorm-gate}
For each integer \(d\ge 1\), define an arithmetic renormalization gate acting on "twisted parameter pairs"
\begin{equation}\label{eq:witt-adams-dirichlet-renorm-gate}
\boxed{\;
\mathsf{Ren}_d:\ (\chi,u)\longmapsto(\chi^d,u^d)\;}
\end{equation}
where \(\chi^d\) is the \(d\)-th power of the character: \(\chi^d(n):=\chi(n)^d\) (still a Dirichlet character), and \(u^d\) represents the Adams operation \(\psi^d\) in the case of coefficient polynomial rings through power replacement (\(\psi^d(u)=u^d\)).
\end{definition}

\begin{proposition}[Renormalization Monad: \texorpdfstring{$(\chi,u)\mapsto(\chi^d,u^d)$}{(chi,u)->(chi^d,u^d)} Semigroup Law and Twist Compatibility]\label{prop:witt-renorm-monoid}
For any \(d,e\ge 1\) there is a semigroup law
\[
\mathsf{Ren}_d\circ\mathsf{Ren}_e=\mathsf{Ren}_{de}.
\]
Moreover, in the calibration of Theorem \ref{thm:witt-dirichlet-twist}, \(\mathsf{Ren}_d\) preserves the form invariance of the twisted--Witt commutative diagram: when \((\chi,u)\) is simultaneously replaced by \((\chi^d,u^d)\), the twisted Möbius formula and the Euler product expression of twisted \(\mathsf{Zeta}_\chi\) still hold term-by-term.
\end{proposition}

\begin{proof}
The semigroup law follows directly from the associativity of powers and the composition law \(\psi^{de}=\psi^d\circ\psi^e\) of Adams operations (see the convention for \(\psi^d\) in Proposition \ref{prop:witt-bridge-functorial}).
Compatibility follows because the formula in Theorem \ref{thm:witt-dirichlet-twist} consists only of the multiplicativity of \(\chi\) and the power replacement property of \(\psi^d\): for any \(d\ge 1\), making the substitution \((\chi,u)\mapsto(\chi^d,u^d)\), all terms still satisfy the same algebraic identity. QED.
\end{proof}

\begin{remark}[Categorified Shadow of Cyclotomic Power Map]\label{rem:witt-renorm-cyclotomic-power-shadow}
On the main text's Dirichlet axis bundle \(G_\infty=\varprojlim_m(\ZZ/m\ZZ)^\times\), the power map \(g\mapsto g^d\) induces \(\chi\mapsto\chi^d\) on the dual side.
Thus the \(\chi\)-direction of \(\mathsf{Ren}_d\) corresponds precisely to the power endomorphism on the cyclotomic quotient, while the \(u\)-direction corresponds to Adams/scale replacement; the two are unified as a single arithmetic rescaling gate within the Big-Witt commutative diagram.
\end{remark}

\begin{proposition}[Necklace Divisibility Law]\label{prop:witt-necklace-div}
For any $n\ge 1$, we have
\[
n\ \Big|\ \sum_{d\mid n}\mu(d)\,a_{n/d}(u^d)\qquad\text{in}\ \mathbb{Z}[u].
\]
\end{proposition}
\begin{proof}
Since $p_n(u)\in\mathbb{Z}[u]$ is the weighted count of primitive orbits, the polynomial is always in the integer coefficient ring; from $n\,p_n(u)=\sum_{d\mid n}\mu(d)a_{n/d}(u^d)$ we obtain the divisibility conclusion. QED.
\end{proof}

\begin{corollary}[$p$-Power Dwork-Type Congruence]\label{cor:witt-dwork-congruence}
Take $n=p^k$ (prime $p$), then
\[
a_{p^k}(u)\equiv a_{p^{k-1}}(u^p)\pmod{p^k}\qquad\text{in}\ \mathbb{Z}[u].
\]
\end{corollary}
\begin{proof}
From $\mu(1)=1,\mu(p)=-1,\mu(p^2)=0$ we get
\(
p^k\,p_{p^k}(u)=a_{p^k}(u)-a_{p^{k-1}}(u^p)
\).
Applying Proposition \ref{prop:witt-necklace-div} immediately gives the conclusion. QED.
\end{proof}

\input{prop__witt-pk-sparsification}

\begin{proposition}[Coefficient Ring Version: Necklace Divisibility Law and Dwork-Type \texorpdfstring{$p$}{p}-Power Congruence]\label{prop:witt-necklace-dwork-R}
Let \(R\) be a commutative \(\ZZ\)-algebra equipped with a family of Adams operations \(\psi^d:R\to R\) (\(d\ge 1\)) satisfying
\(\psi^1=\mathrm{id}\), \(\psi^{de}=\psi^d\circ\psi^e\), and in the case of parameter polynomial rings \(\psi^d\) is given by variable-wise power replacement (e.g., on \(\ZZ[u_1,\dots,u_m]\) it is \(u_i\mapsto u_i^d\)).
Let \(B(\mathbf u)\) be a weighted adjacency matrix family induced by a finite weighted synchronization graph/finite state kernel, with entries in \(R\) (thus each edge carries a weight that is an element of \(R\)), and define
\[
a_n:=\Tr\bigl(B(\mathbf u)^n\bigr)\in R,\qquad n\ge 1.
\]
Define the primitive weighted spectral polynomial as
\[
p_n:=\sum_{\tau\ \mathrm{primitive},\,|\tau|=n} w(\tau)\in R,
\]
where \(w(\tau)\in R\) is the weight product of that primitive orbit (for multiple parameters it is the product of corresponding monomials/weights).
Then for each \(n\ge 1\) there is a strict identity (Witt/Möbius inversion)
\[
n\,p_n=\sum_{d\mid n}\mu(d)\,\psi^d\!\bigl(a_{n/d}\bigr)\in R.
\]
This yields two arithmetic constraints that hold for any such "Euler product kernel":
\begin{enumerate}
  \item \textbf{(Necklace divisibility)} \(\sum_{d\mid n}\mu(d)\,\psi^d(a_{n/d})\in nR\);
  \item \textbf{(Dwork-type congruence)} For any prime \(p\) and \(k\ge 1\), we have
  \[
  a_{p^k}\equiv \psi^{p}\!\bigl(a_{p^{k-1}}\bigr)\pmod{p^k}\qquad\text{in}\ R.
  \]
\end{enumerate}
\end{proposition}
\begin{proof}
In the finite graph/SFT kernel calibration, periodic points (closed loops) and primitive orbits satisfy the standard relation
\[
a_n=\sum_{d\mid n} d\,\psi^{n/d}(p_d)\qquad(n\ge 1),
\]
whose meaning is: loops of length \(n\) decompose by minimal period \(d\), each primitive orbit contributes \(d\) phase starting points under \(n/d\) power repetition; weights under power repetition are acted upon by the Adams operation \(\psi^{n/d}\) (in multiparameter polynomial rings this is variable-wise power replacement \(x\mapsto x^{n/d}\)).
Applying Möbius inversion to the above formula yields
\(
n\,p_n=\sum_{d\mid n}\mu(d)\psi^d(a_{n/d})
\),
and \(p_n\in R\) is guaranteed by its definition as a weighted sum of primitive orbits.
When \(n=p^k\), since \(\mu(p^j)=0\) (\(j\ge 2\)), the above formula collapses to
\(
p^k p_{p^k}=a_{p^k}-\psi^{p}(a_{p^{k-1}})
\),
thus giving the displayed \(\bmod p^k\) congruence. QED.
\end{proof}

\begin{corollary}[\texorpdfstring{$p$}{p}-adic Supercongruence: \texorpdfstring{$a_{p^k}(-1)$}{a_{p^k}(-1)} is Divisible by \texorpdfstring{$p^k$}{p^k}]\label{cor:witt-dwork-supercongruence-minus1}
Take any odd prime \(p\). Then for any \(k\ge 1\) we have
\[
\boxed{\ a_{p^k}(-1)\equiv 0\pmod{p^k}\ }.
\]
\end{corollary}
\begin{proof}
By Corollary \ref{cor:witt-dwork-congruence}, for \(k\ge 1\) we have
\[
a_{p^k}(u)\equiv a_{p^{k-1}}(u^p)\pmod{p^k}\qquad\text{in}\ \ZZ[u].
\]
Specializing at \(u=-1\) and using that \(p\) is odd (so \((-1)^p=-1\)) gives
\[
a_{p^k}(-1)\equiv a_{p^{k-1}}(-1)\pmod{p^k}.
\]
Since \(p^{k-1}\) is odd, by the conclusion \(a_n(-1)=0\) of Proposition \ref{prop:trace-palindrome} (odd \(n\)) we know \(a_{p^{k-1}}(-1)=0\).
Thus \(a_{p^k}(-1)\equiv 0\pmod{p^k}\). QED.
\end{proof}

\begin{remark}[\texorpdfstring{$p$}{p}-Power Direction Witt Interpretation: Frobenius-Compatible Ghost Components]\label{rem:witt-ptypical-ghost}
Viewing \(a_n(u)\) as ghost coordinates on the Big-Witt (necklace) ring, then \(\{a_{p^k}(u)\}_{k\ge 0}\) is its ghost subsequence on \(p\)-power indices.
The Dwork-type congruence in Corollary \ref{cor:witt-dwork-congruence} shows that this subsequence satisfies strong \(p\)-adic compatibility under the Adams/Frobenius operation \(u\mapsto u^p\):
\[
a_{p^k}(u)-a_{p^{k-1}}(u^p)\in p^k\,\ZZ[u].
\]
Thus the primitive/trace bridge package of the synchronization kernel is auditable not only on the complex analytic side (\(\zeta\), Euler product), but also carries a natural "Frobenius-lift" constraint on the \(p\)-adic integrality side.
In future considerations of \(p\)-adic direction interpolation or comparison of "arithmetic fingerprints" of different kernels, this congruence can serve as the most basic auditable input.
\end{remark}

\begin{proposition}[\texorpdfstring{$p$}{p}-Power Primitive Coordinate Explicit Formula]\label{prop:witt-ppower-primitive-coordinate}
For any prime \(p\) and \(k\ge 1\), the Big-Witt/Möbius inversion at \(n=p^k\) collapses to
\[
\boxed{\;
p_{p^k}(u)=\frac{a_{p^k}(u)-a_{p^{k-1}}(u^p)}{p^k}\in\ZZ[u]\; }.
\]
Thus Corollary \ref{cor:witt-dwork-congruence} is equivalent to "\(p\)-power direction primitive coordinate \(p_{p^k}(u)\) is well-defined in the integer coefficient ring," and provides an auditable \(p\)-adic difference fingerprint sequence \(\{p_{p^k}(u)\}_{k\ge 0}\).
\end{proposition}
\begin{proof}
By definition
\(
p_n(u)=\frac{1}{n}\sum_{d\mid n}\mu(d)\,a_{n/d}(u^d)
\).
When \(n=p^k\), divisors are only \(p^j\), and \(\mu(p^j)=0\) (\(j\ge 2\)), so only \(d=1\) and \(d=p\) remain, yielding
\[
p_{p^k}(u)=\frac{1}{p^k}\bigl(a_{p^k}(u)-a_{p^{k-1}}(u^p)\bigr).
\]
By Proposition \ref{prop:witt-necklace-div} the right side is indeed a \(\ZZ[u]\) element. QED.
\end{proof}

\begin{lemma}[Kernel-Version Dieudonné--Dwork Integrality Criterion (Formal Power Series)]\label{lem:witt-dieudonne-dwork}
Let \(R\) be a \(p\)-torsion-free \(p\)-adically complete ring (e.g.\ \(R=\ZZ_p[u]\)), and define the Frobenius lift operator
\[
\varphi:\ R[[T]]\to R[[T]],\qquad
\varphi\!\left(\sum_{n\ge 0}c_n(u)T^n\right):=\sum_{n\ge 0}c_n(u^p)\,T^{pn}.
\]
If \(F(T)\in 1+T\,R\otimes_{\ZZ_p}\QQ_p[[T]]\) satisfies
\[
\boxed{\;
\frac{\varphi(F(T))}{F(T)^p}\in 1+pT\,R[[T]]\;},
\]
then necessarily \(F(T)\in 1+T\,R[[T]]\) (i.e., all coefficients of \(F\) are \(p\)-integral).
\end{lemma}
\begin{proof}
This is a commonly used integrality criterion in \(p\)-typical Witt/Artin--Hasse theory (Dieudonné--Dwork form).
In the form needed for this paper, one can prove it as follows: take the \(p\)-adic logarithm \(\log F\) of \(F\) (which converges in the \(T\)-adic topology), rewrite the condition as
\(
\varphi(\log F)-p\log F\in pT\,R[[T]]
\);
then by induction on the degree of \(T\), one can successively deduce that the coefficients of \(\log F\) are all in \(R\), so \(F=\exp(\log F)\) is also in \(1+T\,R[[T]]\). QED.
\end{proof}

\begin{theorem}[Kernel-Version Dwork Lemma \texorpdfstring{$p$}{p}-Typical Integrality: Local \texorpdfstring{$p$}{p}-Factor Closed in \texorpdfstring{$\ZZ_p[u]$}{Z\_p[u]}]\label{thm:witt-dwork-zp-integrality}
Fix a prime \(p\). Define
\[
S_p(T;u):=\sum_{k\ge 0}\frac{a_{p^k}(u)}{p^k}\,T^{p^k}\in \QQ_p[u][[T]],
\qquad
Z_p(T;u):=\exp\!\bigl(S_p(T;u)\bigr).
\]
Then we have
\[
\boxed{\;
Z_p(T;u)\in 1+T\,\ZZ_p[u][[T]]\; }.
\]
\end{theorem}
\begin{proof}
Take \(R=\ZZ_p[u]\) in Lemma \ref{lem:witt-dieudonne-dwork}.
Applying \(\varphi\) to \(S_p\) gives
\[
\varphi(S_p)(T;u)
=\sum_{k\ge 0}\frac{a_{p^k}(u^p)}{p^k}\,T^{p^{k+1}}
=\sum_{k\ge 1}\frac{a_{p^{k-1}}(u^p)}{p^{k-1}}\,T^{p^{k}}.
\]
On the other hand
\[
pS_p(T;u)
=p\,a_1(u)\,T+\sum_{k\ge 1}\frac{a_{p^k}(u)}{p^{k-1}}\,T^{p^k}.
\]
Therefore
\[
\varphi(S_p)(T;u)-pS_p(T;u)
=-p\,a_1(u)\,T+\sum_{k\ge 1}\frac{a_{p^{k-1}}(u^p)-a_{p^k}(u)}{p^{k-1}}\,T^{p^k}.
\]
By Corollary \ref{cor:witt-dwork-congruence}, for \(k\ge 1\) we have
\(
a_{p^k}(u)-a_{p^{k-1}}(u^p)\in p^k\ZZ[u]
\),
so each term
\(
\frac{a_{p^{k-1}}(u^p)-a_{p^k}(u)}{p^{k-1}}
\)
is in \(p\,\ZZ_p[u]\). Also \(-p a_1(u)T\in pT\,\ZZ_p[u][[T]]\), thus
\[
\varphi(S_p)-pS_p\in pT\,\ZZ_p[u][[T]].
\]
Let \(Z_p=\exp(S_p)\), then
\[
\frac{\varphi(Z_p)}{Z_p^p}
=\exp\!\bigl(\varphi(S_p)-pS_p\bigr)\in 1+pT\,\ZZ_p[u][[T]].
\]
By Lemma \ref{lem:witt-dieudonne-dwork} we get \(Z_p\in 1+T\,\ZZ_p[u][[T]]\). QED.
\end{proof}

\begin{proposition}[Unit-Root Limit Invariant: \texorpdfstring{$u^p=u$}{u^p=u} Implies \texorpdfstring{$a_{p^k}(u)$}{a_{p^k}(u)} Converges in \texorpdfstring{$\ZZ_p$}{Z_p}]\label{prop:witt-unit-root-limit}
Let \(p\) be an odd prime. If \(u\) satisfies \(u^p=u\) (e.g., \(u\in\mu_{p-1}\subset\ZZ_p^\times\) is a Teichm\"uller unit), then the sequence \(\{a_{p^k}(u)\}_{k\ge 0}\) is Cauchy in \(\ZZ_p\), and converges to the limit
\[
\boxed{\;
A_p(u):=\lim_{k\to\infty}a_{p^k}(u)\in \ZZ_p\; }.
\]
Moreover, there is an explicit \(p\)-adic error bound
\[
v_p\!\bigl(a_{p^k}(u)-a_{p^{k-1}}(u)\bigr)\ge k\qquad(k\ge 1).
\]
\end{proposition}
\begin{proof}
Under the condition \(u^p=u\), Corollary \ref{cor:witt-dwork-congruence} degenerates to
\(
a_{p^k}(u)\equiv a_{p^{k-1}}(u)\ (\mathrm{mod}\ p^k)
\),
i.e., \(a_{p^k}(u)-a_{p^{k-1}}(u)\in p^k\ZZ_p\), thus giving the displayed \(p\)-adic valuation bound.
Therefore for any \(r>s\), we have
\[
a_{p^r}(u)-a_{p^s}(u)=\sum_{k=s+1}^{r}\bigl(a_{p^k}(u)-a_{p^{k-1}}(u)\bigr)\in p^{s+1}\ZZ_p,
\]
so it is Cauchy and converges in \(\ZZ_p\). QED.
\end{proof}

\begin{remark}[Unit-Root Limit and Self-Dual Compatibility]\label{rem:witt-unit-root-self-dual}
If we further use the palindrome relation \(a_n(u)=u^n a_n(1/u)\) of Proposition \ref{prop:trace-palindrome}, then in the Teichm\"uller case \(u^p=u\) we have \(u^{p^k}=u\), thus
\[
a_{p^k}(u)=u\,a_{p^k}(1/u).
\]
Letting \(k\to\infty\) and using the convergence of Proposition \ref{prop:witt-unit-root-limit}, we can obtain that the limit invariant satisfies
\[
\boxed{\ A_p(u)=u\,A_p(1/u)\ }.
\]
This provides an auditable constraint for the unit-root fingerprint under self-dual symmetry.
\end{remark}

\subsubsection{Self-Dual Symmetry of the Synchronization Kernel}\label{app:kernel-self-dual}
For the synchronization kernel there exists a state involution $\iota$, whose action is
\[
\iota:\quad
000\leftrightarrow 101,\quad
001\leftrightarrow 100,\quad
002\leftrightarrow 010,\quad
0\overline{1}2\leftrightarrow 11\overline{1},\quad
1\overline{1}2\leftrightarrow 01\overline{1}.
\]
Let $\Pi$ be the matrix of this permutation, then
\[
\Pi^{-1}B_0\Pi=B_1,\qquad \Pi^{-1}B_1\Pi=B_0,
\]
thus
\[
\Pi^{-1}B(u)\Pi=uB(1/u).
\]
\begin{lemma}[\texorpdfstring{$\Pi$}{Pi}-Even/Odd Component and Unique Decomposition of \texorpdfstring{$B(u)$}{B(u)}]\label{lem:sync-kernel-even-odd-splitting}
Define
\[
B_{\mathrm{sym}}:=\frac{B_0+B_1}{2},\qquad
B_{\mathrm{asym}}:=\frac{B_1-B_0}{2}.
\]
Then for all \(u\in\CC\) we have
\[
\boxed{\ B(u)=(1+u)\,B_{\mathrm{sym}}+(u-1)\,B_{\mathrm{asym}}\ }.
\]
Moreover, under the self-dual commutation relation \(\Pi^{-1}B_0\Pi=B_1,\ \Pi^{-1}B_1\Pi=B_0\), we have
\[
\boxed{\ \Pi^{-1}B_{\mathrm{sym}}\Pi=B_{\mathrm{sym}},\qquad
\Pi^{-1}B_{\mathrm{asym}}\Pi=-B_{\mathrm{asym}}\ }.
\]
Furthermore, the above \((1+u)\)--\((u-1)\) decomposition is unique in the \(\QQ\)-linear sense: if there exist \(S,A\) such that
\(
B(u)=(1+u)S+(u-1)A
\), then necessarily \(S=B_{\mathrm{sym}}\) and \(A=B_{\mathrm{asym}}\).
\end{lemma}

\begin{proof}
By definition, expanding directly
\[
(1+u)B_{\mathrm{sym}}+(u-1)B_{\mathrm{asym}}
=\frac12\Bigl((1+u)(B_0+B_1)+(u-1)(B_1-B_0)\Bigr)
=B_0+uB_1
=B(u).
\]
Then from \(\Pi\) exchanging \(B_0,B_1\) we get
\[
\Pi^{-1}B_{\mathrm{sym}}\Pi
=\frac{1}{2}\bigl(\Pi^{-1}B_0\Pi+\Pi^{-1}B_1\Pi\bigr)
=\frac{B_1+B_0}{2}
=B_{\mathrm{sym}},
\]
\[
\Pi^{-1}B_{\mathrm{asym}}\Pi
=\frac{1}{2}\bigl(\Pi^{-1}B_1\Pi-\Pi^{-1}B_0\Pi\bigr)
=\frac{B_0-B_1}{2}
=-B_{\mathrm{asym}}.
\]
Uniqueness is given by comparing at the two points \(u=0,1\): if \(B(u)=(1+u)S+(u-1)A\), then
\(
B(0)=S-A=B_0
\)
and
\(
B(1)=2S=B_0+B_1
\)
jointly solve to give \(S=(B_0+B_1)/2=B_{\mathrm{sym}}\), and substituting back gives \(A=(B_1-B_0)/2=B_{\mathrm{asym}}\).
QED.
\end{proof}

\begin{remark}[Computable Bridge Package: Matrix-Level Verification Checklist for Lemma \ref{lem:sync-kernel-even-odd-splitting}]\label{rem:sync-kernel-even-odd-bridge}
Given integer matrices \((B_0,B_1)\) and permutation matrix \(\Pi\), verify the following equations one by one:
\begin{enumerate}
  \item \textbf{Self-dual commutation:} Verify \(\Pi^{-1}B_0\Pi=B_1\) and \(\Pi^{-1}B_1\Pi=B_0\) (under permutation matrix equivalently \(\Pi^\top B_0\Pi=B_1\)).
  \item \textbf{Even/odd components:} Compute \(2B_{\mathrm{sym}}=B_0+B_1\), \(2B_{\mathrm{asym}}=B_1-B_0\), and verify
  \[
  \Pi^{-1}(2B_{\mathrm{sym}})\Pi=2B_{\mathrm{sym}},\qquad
  \Pi^{-1}(2B_{\mathrm{asym}})\Pi=-2B_{\mathrm{asym}}.
  \]
  \item \textbf{\(u\)-dependent reconstruction:} For any sampled \(u\) (e.g., random complex number), verify
  \(
  B(u)=B_0+uB_1=(1+u)B_{\mathrm{sym}}+(u-1)B_{\mathrm{asym}}
  \).
\end{enumerate}
In the basis of the \(\Pi\) eigende composition of Proposition \ref{prop:self-dual-normal-form-1pmu}, the above components correspond to
\(
B_{\mathrm{sym}}=\bigl(\begin{smallmatrix}X&0\\0&W\end{smallmatrix}\bigr)
\)
and
\(
B_{\mathrm{asym}}=\bigl(\begin{smallmatrix}0&-Y\\-Z&0\end{smallmatrix}\bigr)
\).
The corresponding one-click audit script is \texttt{scripts/exp\_sync\_kernel\_weighted\_self\_dual\_bridge.py}.
\end{remark}

\begin{definition}[Completed Matrix]\label{def:completed-matrix}
Let \(u=e^{\theta}\) (\(\theta\in\CC\)), and define the completed matrix
\[
\widetilde B(\theta):=e^{-\theta/2}B(e^{\theta}).
\]
\end{definition}

\begin{lemma}[Self-Dual Similarity of Completion]\label{lem:completed-similarity}
Under the above notation, for all \(\theta\) we have
\[
\boxed{\ \Pi^{-1}\widetilde B(\theta)\Pi=\widetilde B(-\theta)\ }.
\]
Thus \(\widetilde B(\theta)\) is similar to \(\widetilde B(-\theta)\); in particular,
\(\det(\lambda I-\widetilde B(\theta))\) as an analytic function of \(\theta\) is an even function.
\end{lemma}
\begin{proof}
By self-duality \(\Pi^{-1}B(u)\Pi=uB(1/u)\), substituting \(u=e^{\theta}\) gives
\(
\Pi^{-1}B(e^\theta)\Pi=e^\theta B(e^{-\theta})
\).
Multiplying both sides by \(e^{-\theta/2}\) immediately yields
\(
\Pi^{-1}\widetilde B(\theta)\Pi=\widetilde B(-\theta)
\).
Similarity implies the characteristic polynomials are equal, so \(\det(\lambda I-\widetilde B(\theta))=\det(\lambda I-\widetilde B(-\theta))\), which is an even function of \(\theta\). QED.
\end{proof}

\begin{lemma}[Ward Identity: \texorpdfstring{$B_{\mathrm{asym}}$}{B_asym} Orthogonal to All Time Traces at First Order]\label{lem:sync-kernel-ward-traces}
Following the notation of Lemma \ref{lem:sync-kernel-even-odd-splitting}, and let
\(
\widetilde B(\theta)=e^{-\theta/2}B(e^{\theta})
\).
Then for any integer \(n\ge 1\) there is a strict identity
\[
\boxed{\ \Tr\!\bigl(B(1)^{n-1}\,B_{\mathrm{asym}}\bigr)=0\ }.
\]
\end{lemma}

\begin{proof}
By Lemma \ref{lem:completed-similarity}, \(\widetilde B(\theta)\) is similar to \(\widetilde B(-\theta)\), so
\(
\Tr(\widetilde B(\theta)^n)
\)
is an even function of \(\theta\), hence its first derivative at \(\theta=0\) is zero:
\[
0=\left.\frac{d}{d\theta}\Tr(\widetilde B(\theta)^n)\right|_{\theta=0}.
\]
On the other hand, for an analytic matrix family \(A(\theta)\) we have
\(
\frac{d}{d\theta}\Tr(A(\theta)^n)
=\sum_{j=0}^{n-1}\Tr\!\bigl(A(\theta)^j A'(\theta)A(\theta)^{n-1-j}\bigr)
=n\,\Tr\!\bigl(A(\theta)^{n-1}A'(\theta)\bigr)
\)
(last step using trace cyclicity).
Taking \(A(\theta)=\widetilde B(\theta)\) and substituting at \(\theta=0\) gives
\[
0=n\,\Tr\!\bigl(\widetilde B(0)^{n-1}\widetilde B'(0)\bigr).
\]
Note that \(\widetilde B(0)=B(1)\), and from
\(
\widetilde B(\theta)=e^{-\theta/2}(B_0+e^\theta B_1)
\)
we directly compute
\[
\widetilde B'(0)=-\frac12(B_0+B_1)+B_1=\frac{B_1-B_0}{2}=B_{\mathrm{asym}}.
\]
Thus we obtain the displayed identity. QED.
\end{proof}

\begin{remark}[Computable Bridge Package: Trace Verification of Lemma \ref{lem:sync-kernel-ward-traces}]\label{rem:sync-kernel-ward-traces-bridge}
For given \((B_0,B_1,\Pi)\), first compute \(B_{\mathrm{asym}}\) and \(B(1)=B_0+B_1\) via Lemma \ref{lem:sync-kernel-even-odd-splitting}, then verify term-by-term for a range of length upper bound \(1\le n\le N\) (e.g., \(N=30\))
\[
\Tr\!\bigl(B(1)^{n-1}B_{\mathrm{asym}}\bigr)=0.
\]
This identity is an exact equation at the integer (or rational) level, thus suitable for independent auditing (not relying on floating-point spectral ordering). The one-click verification script is \texttt{scripts/exp\_sync\_kernel\_weighted\_self\_dual\_bridge.py}.
\end{remark}

\begin{theorem}[Spectral Tangent Lock: All Spectral Logarithmic Slopes at \texorpdfstring{$u=1$}{u=1} Are Identically \texorpdfstring{$1/2$}{1/2}]\label{thm:sync-spectral-tangent-lock}
Let \(\widetilde B(\theta)\) be as in Definition \ref{def:completed-matrix}.
Take any \textbf{simple nonzero eigenvalue} \(\mu_0\neq 0\) of \(\widetilde B(0)=B(1)\) (i.e., it is a simple root in the characteristic polynomial).
Then there exists a unique analytic function \(\mu(\theta)\) (in some neighborhood of \(\theta=0\)) satisfying
\[
\mu(0)=\mu_0,\qquad \mu(\theta)\in\mathrm{spec}\bigl(\widetilde B(\theta)\bigr),
\]
and this branch satisfies strict evenness
\[
\boxed{\ \mu(\theta)=\mu(-\theta)\ } \qquad(\text{thus } \mu'(0)=0).
\]
Let
\[
\lambda(\theta):=e^{\theta/2}\mu(\theta),
\]
then \(\lambda(\theta)\) is an analytic eigenvalue branch of matrix \(B(e^{\theta})\) (and \(\lambda(0)=\mu_0\neq 0\)), satisfying
\[
\boxed{\ \left.\frac{\mathrm{d}}{\mathrm{d}\theta}\log\lambda(\theta)\right|_{\theta=0}=\frac12\ }.
\]
\end{theorem}
\begin{proof}
By Lemma \ref{lem:completed-similarity}, \(\chi(\lambda,\theta):=\det(\lambda I-\widetilde B(\theta))\) is an even function of \(\theta\), and \(\chi(\mu_0,0)=0\).
Since \(\mu_0\) is a simple eigenvalue, we have \(\partial_\lambda\chi(\mu_0,0)\neq 0\), so by the implicit function theorem there exists a unique analytic branch \(\mu(\theta)\) such that \(\chi(\mu(\theta),\theta)=0\) and \(\mu(0)=\mu_0\).
On the other hand, evenness gives \(\chi(\lambda,\theta)=\chi(\lambda,-\theta)\), so \(\mu(-\theta)\) also satisfies the same equation and takes the same initial value \(\mu(-0)=\mu_0\); by uniqueness we obtain \(\mu(\theta)=\mu(-\theta)\), hence \(\mu'(0)=0\).
Finally, from \(B(e^\theta)=e^{\theta/2}\widetilde B(\theta)\) we know \(\lambda(\theta)=e^{\theta/2}\mu(\theta)\) is an eigenvalue branch of \(B(e^\theta)\), and
\[
\left.\frac{\mathrm{d}}{\mathrm{d}\theta}\log\lambda(\theta)\right|_{\theta=0}
=\left.\frac{\mathrm{d}}{\mathrm{d}\theta}\left(\frac{\theta}{2}+\log\mu(\theta)\right)\right|_{\theta=0}
=\frac12+\frac{\mu'(0)}{\mu(0)}
=\frac12.
\]
QED.
\end{proof}

\begin{corollary}[The 6 Nonzero Spectral Branches of the Synchronization Kernel Simultaneously Lock at \texorpdfstring{$u=1$}{u=1}]\label{cor:sync-spectral-tangent-lock-6}
For this synchronization kernel, matrix \(B(1)\) has exactly \(6\) nonzero spectral branches, all simple eigenvalues at \(u=1\) (its characteristic polynomial is degree \(6\) of \(\det(I-zB(1))\), see Appendix \ref{app:sync-kernel-basic} for explicit factorization).
Thus all these \(6\) nonzero spectral branches satisfy the conclusion of Theorem \ref{thm:sync-spectral-tangent-lock}:
under the parameter \(u=e^\theta\), \(\left.\frac{\mathrm{d}}{\mathrm{d}\theta}\log\lambda_j(\theta)\right|_{\theta=0}=\frac12\).
\end{corollary}

Thus
\[
\Delta(z,u)=\det(I-zB(u))=\det(I-uzB(1/u))=\Delta(uz,1/u).
\]

\begin{proposition}[Self-Dual-Induced \texorpdfstring{$(1+u)$}{(1+u)}--\texorpdfstring{$(1-u)$}{(1-u)} Two-Parameter Normal Form]\label{prop:self-dual-normal-form-1pmu}
Suppose \(B(u)=B_0+uB_1\) and the self-dual permutation matrix \(\Pi\) satisfies \(\Pi^{-1}B_0\Pi=B_1\).
Let
\[
V_+:=\ker(\Pi-I),\qquad V_-:=\ker(\Pi+I),
\qquad V=V_+\oplus V_-,
\]
and write \(B_0\) in the decomposition \(V_+\oplus V_-\) as block operator
\[
B_0=
\begin{pmatrix}
X & Y\\
Z & W
\end{pmatrix}.
\]
Then for all \(u\in\CC\) we have
\[
\boxed{
B(u)=B_0+u\Pi^{-1}B_0\Pi
=
\begin{pmatrix}
(1+u)X & (1-u)Y\\
(1-u)Z & (1+u)W
\end{pmatrix}.
}
\]
For this synchronization kernel, \(\Pi\) is a direct sum of five interchanges, so \(\dim V_+=\dim V_-=5\).
\end{proposition}
\begin{proof}
In the decomposition \(V=V_+\oplus V_-\), \(\Pi\) acts as \(\Pi|_{V_+}=+I,\ \Pi|_{V_-}=-I\), so for any block operator \(T=\bigl(\begin{smallmatrix}X&Y\\ Z&W\end{smallmatrix}\bigr)\) we have
\(
\Pi^{-1}T\Pi=\bigl(\begin{smallmatrix}X&-Y\\ -Z&W\end{smallmatrix}\bigr)
\).
Taking \(T=B_0\) and substituting into \(B(u)=B_0+u\Pi^{-1}B_0\Pi\) immediately gives the conclusion. QED.
\end{proof}

\begin{corollary}[\texorpdfstring{$u=1$}{u=1} Forced Decoupling]\label{cor:self-dual-u1-decouple}
In the block representation of Proposition \ref{prop:self-dual-normal-form-1pmu},
\[
\boxed{\ B(1)=
\begin{pmatrix}
2X & 0\\
0 & 2W
\end{pmatrix}\ }.
\]
Thus cross-subspace coupling terms appear only through the coefficient \((1-u)\) and are forcibly suppressed at \(u=1\).
\end{corollary}

\begin{corollary}[\texorpdfstring{$u=1$}{u=1} Two-Channel \texorpdfstring{$\Delta/\zeta$}{Delta/zeta} Factorization]\label{cor:self-dual-u1-two-channel-zeta}
Following the block notation of Proposition \ref{prop:self-dual-normal-form-1pmu}, at \(u=1\) there is a strict factor decomposition
\[
\boxed{\ \Delta(z,1)=\det(I-zB(1))=\det(I-2zX)\cdot\det(I-2zW)\ }.
\]
Thus defining two "channel \(\zeta\)" as
\[
\zeta_{V_+}(z):=\det(I-2zX)^{-1},\qquad
\zeta_{V_-}(z):=\det(I-2zW)^{-1},
\]
we have
\[
\boxed{\ \zeta(z,1)=\Delta(z,1)^{-1}=\zeta_{V_+}(z)\,\zeta_{V_-}(z)\ }.
\]
\end{corollary}

\begin{proof}
By Corollary \ref{cor:self-dual-u1-decouple}, \(B(1)\) in the decomposition \(V_+\oplus V_-\) is block diagonal \(\mathrm{diag}(2X,2W)\).
Taking \(\det(I-z\cdot)\) of a block diagonal matrix immediately gives
\(
\det(I-zB(1))=\det(I-2zX)\det(I-2zW)
\).
The \(\zeta\) factorization is obtained by taking reciprocals. QED.
\end{proof}

\begin{remark}[Computable Bridge Package: Block Verification of Corollary \ref{cor:self-dual-u1-two-channel-zeta}]\label{rem:sync-kernel-u1-two-channel-bridge}
For given \((B_0,B_1,\Pi)\), first diagonalize \(\Pi\) (take any basis of \(V_+\oplus V_-\), equivalently pair states by \(\Pi\)'s \(2\)-cycles and take \(\pm\) combinations), and accordingly write \(B(1)\) as block diagonal \(\mathrm{diag}(2X,2W)\).
Then:
\begin{enumerate}
  \item \textbf{Matrix level:} Verify the off-diagonal blocks of \(B(1)\) are zero;
  \item \textbf{Determinant level:} For a set of sampled \(z\) verify \(\det(I-zB(1))=\det(I-2zX)\det(I-2zW)\);
  \item \textbf{Channel \(\zeta\):} Compute \(\zeta_{V_\pm}\) separately and verify \(\zeta(z,1)=\zeta_{V_+}(z)\zeta_{V_-}(z)\).
\end{enumerate}
The one-click verification script is \texttt{scripts/exp\_sync\_kernel\_weighted\_self\_dual\_bridge.py}.
\end{remark}

\begin{corollary}[\texorpdfstring{$\Delta(z,u)$}{Delta(z,u)} Dependence Compression in \texorpdfstring{$(1-u)^2$}{(1-u)^2}]\label{cor:self-dual-delta-even-in-1mu}
In the normal form of Proposition \ref{prop:self-dual-normal-form-1pmu}, for each fixed \(z\),
\(\Delta(z,u)=\det(I-zB(u))\)
as a polynomial in \(u\), appears only in even powers of \((1-u)\); equivalently,
\[
\boxed{\ \Delta(z,u)\in \ZZ[z]\bigl[(1+u),\ (1-u)^2\bigr]\ }.
\]
Similarly, expanding \(a_n(u)=\Tr(B(u)^n)\) as a \(u\)-polynomial, its coefficients can also be rewritten as elements in \(\ZZ[(1+u),(1-u)^2]\).
\end{corollary}
\begin{proof}
The permutation expansion of the determinant consists of permutation cycles; each cycle must return to the starting point, so it crosses \(V_+\leftrightarrow V_-\) an even number of times.
Crossings come only from off-diagonal blocks \((1-u)Y,(1-u)Z\), so in each term the total index of \((1-u)\) is even, thus \(\Delta\) is a polynomial in \((1-u)^2\).
The same holds for \(\Tr(B(u)^n)\): it is a sum of weights of closed paths of length \(n\), and each closed path has an even number of subspace flips.
QED.
\end{proof}

Moreover, \(\Delta\) for this kernel can be directly recovered as an explicit polynomial from the sextic equation \(F(\lambda,u)=0\) (and pointwise verified consistent with the \(10\times 10\) matrix definition):
\input{sections/generated/eq_sync_kernel_weighted_delta_explicit}
\begin{proposition}[Completed Determinant Bivariate Closed Form: \texorpdfstring{$\Delta(z,u)=\widehat\Delta(w,s)$}{Delta=hatDelta}]\label{prop:sync-kernel-weighted-hat-delta}
Take completed coordinates
\[
u=r^2,\qquad w:=zr=z\sqrt{u},\qquad s:=r+r^{-1}=\sqrt{u}+\frac{1}{\sqrt{u}}.
\]
Then there exists an integer coefficient bivariate polynomial \(\widehat\Delta(w,s)\in\ZZ[w,s]\) such that
\[
\boxed{\ \Delta(z,u)=\widehat\Delta\!\left(z\sqrt{u},\ \sqrt{u}+\frac{1}{\sqrt{u}}\right)\ }.
\]
For this kernel, \(\widehat\Delta\) can be taken as the closed form of equation \eqref{eq:sync-kernel-weighted-delta-completed} (generated directly by script and verified term-by-term consistent with \eqref{eq:sync-kernel-weighted-delta-explicit}):
\input{sections/generated/eq_sync_kernel_weighted_delta_completed}
\end{proposition}
\begin{proof}
By self-duality \(\Delta(z,u)=\Delta(uz,1/u)\), writing \(u=r^2\) and letting \(w=zr\), we know
\(
\Delta(w/r,r^2)
\)
is invariant under \(r\mapsto r^{-1}\).
The invariant subring identity
\(
\ZZ[r,r^{-1}]^{\,r\mapsto r^{-1}}=\ZZ[r+r^{-1}]=\ZZ[s]
\)
ensures it can be written as some \(\widehat\Delta(w,s)\in\ZZ[w,s]\).
Substituting \eqref{eq:sync-kernel-weighted-delta-explicit} into \(u=r^2,z=w/r\) and simplifying on the invariant basis \(C_k(s)=r^k+r^{-k}\) yields \eqref{eq:sync-kernel-weighted-delta-completed}; the script \texttt{scripts/exp\_sync\_kernel\_weighted\_delta\_completed.py} performs symbolic verification of this identity. QED.
\end{proof}
